% Generated by scripts/build-paper.py — edit paper/src/survey.src.md
\section{Lineage and an analytical frame}\label{sec:lineage-and-an-analytical-frame}

\subsection{Lineages}

Typed decisions recombine long-standing ideas. \textbf{Discriminative classification} with bidirectional encoders \citep{arxiv1810_04805,arxiv2412_13663}, few-shot classifiers \citep{arxiv2209_11055} and generalist lightweight classifiers \citep{arxiv2508_07662,arxiv2311_08526} are the natural baselines, and zero-shot labelling via label descriptions anticipated runtime-defined options \citep{arxiv1909_00161}. Task-agnostic discriminative modules \citep{arxiv2306_07536} and domain probabilistic decision systems such as SalesRLAgent \citep{arxiv2503_23303} predate the product. They establish context, not shared architecture or priority. \textbf{Calibration} and temperature scaling \citep{arxiv1706_04599} and the failure of calibration under shift \citep{arxiv1906_02530} frame every probability claim, alongside the gap between models' self-knowledge and verbalised confidence \citep{arxiv2207_05221}, semantic uncertainty \citep{arxiv2302_09664} and RL-trained confidence expression \citep{arxiv2503_02623}. \textbf{Selective prediction} \citep{arxiv1705_08500,arxiv1901_09192}, learning to defer \citep{arxiv2006_01862} and conformal methods \citep{arxiv2107_07511,arxiv2310_05921} supply the theory of acting on confidence. Decision-aware calibration measures what matters downstream \citep{arxiv2404_13503,arxiv2504_15582,arxiv2511_13699,arxiv2510_07750}. \textbf{Structured output} research shows that constrained decoding guarantees validity but not semantics \citep{arxiv2109_05093,arxiv2307_09702,arxiv2411_15100}, that prefix sharing and scheduling yield their own speed-ups \citep{arxiv2312_07104}, and that schema compliance and forced forms need their own benchmarks \citep{arxiv2501_10868,arxiv2604_25359,arxiv2607_20492}. \textbf{Judges and reward models} \citep{arxiv2306_05685,arxiv2310_08491,arxiv2403_13787}, \textbf{routing and cascades} \citep{arxiv2305_05176,arxiv2406_18665,arxiv2601_07206,arxiv2510_01237} and \textbf{label and format sensitivity} \citep{arxiv2309_04992,arxiv2310_11324,arxiv2608_08254} complete the background.

\subsection{Five places a claim can live}

We read each study by where its claim sits on the path from a declared question to a workflow outcome (Table 1). This is an analytical organisation for comparing evidence. It is not the internal architecture of any model, and we do not claim it as the first such taxonomy.

\begin{table}[tbp]
\centering\footnotesize
\caption{Five places a claim can live: an analytical organisation, not a model architecture.}\label{tab:stages}
\begin{tabularx}{\linewidth}{p{0.16\linewidth}XXr}
\toprule
\textbf{Stage} & \textbf{Question} & \textbf{Evidence to look for} & \textbf{Studies} \\
\midrule
Decision contract & What exactly is being asked, and what may the answer be? & Type validity; name and order invariance; closed-set failure when the right option is missing. & 6 \\
Inference \& readout & How is a distribution over the declared options produced? & Same-backbone ablations; readout geometry; what is and is not disclosed. & 5 \\
Probability \& calibration & What does a returned probability or confidence mean, and for which population? & Per-task reliability, high-confidence errors, recalibration on held-out labels. & 7 \\
Selective control & When should software act, ask, escalate or abstain? & Coverage at fixed risk; escalation rate; thresholds chosen on a separate selection set. & 6 \\
Downstream use & Does the whole workflow get better, cheaper or faster? & System-level baselines; module attribution; matched measurement scope. & 16 \\
\bottomrule
\end{tabularx}
\end{table}
